\documentclass[11pt]{article}
% Shared preamble for the bite-sized LaTeX doc series (docs/latex/*.tex).
% \input this at the top of each numbered file rather than duplicating the
% package list -- keeps every file's own content close to (and comfortably
% under) 10 pages.
\usepackage[a4paper,margin=2.5cm]{geometry}
\usepackage{amsmath,amssymb}
\usepackage{hyperref}
\usepackage[backend=biber,style=numeric,sorting=none]{biblatex}
\usepackage{listings}
\usepackage{xcolor}
\usepackage{enumitem}
\addbibresource{references.bib}

\lstdefinelanguage{Rust}{
  keywords={fn,let,mut,pub,struct,enum,impl,match,if,else,for,while,return,use,mod,self,Self,const,static},
  keywordstyle=\color{blue}\bfseries,
  comment=[l]{//},
  commentstyle=\color{gray}\itshape,
  string=[b]",
  stringstyle=\color{teal},
  basicstyle=\ttfamily\small,
  showstringspaces=false,
  breaklines=true,
}
\lstset{language=Rust}

\hypersetup{colorlinks=true, linkcolor=blue, citecolor=blue, urlcolor=blue}


\title{The Non-Analytic Critical-Region Term \\ \large Part 4: Function-by-Function Walkthrough}
\author{outram-park-fork-coolprop}
\date{2026-07-10}

\begin{document}
\maketitle

\begin{abstract}
This document traces one concrete call — computing Water's pressure at its
own critical point — through every function it touches, from the public API
down to the non-analytic term evaluator from Part 3, and shows the actual
verified numerical result. Read alongside the source files named in each
section; nothing here should require jumping to a file not named.
\end{abstract}

\section{The entry point: \texttt{state\_trho}}

Everything starts from \texttt{props::state\_trho(fluid, t, rho)}
(\texttt{src/props.rs}) — the crate's most basic query, ``give me every
property of this fluid at this temperature and density'':
\begin{lstlisting}[language=Rust]
let eos: &FluidEos = fluid.eos();
let delta = rho_molar / eos.rho_reducing;
let tau = eos.t_reducing / t;
let res = eos.residual_derivs(delta, tau);
let ide = eos.ideal_derivs(delta, tau);
// ... combine res, ide into p, u, h, s, cv, cp, speed_of_sound
\end{lstlisting}
For our example, \texttt{fluid = Fluid::Water}, and we choose \texttt{t}
and \texttt{rho} to land \emph{exactly} on Water's own critical point:
\texttt{t = 647.096} (Kelvin) and \texttt{rho} = the critical point's mass
density, $322.0\ \mathrm{kg/m^3}$ (obtained from
\texttt{eos.rho\_critical * eos.molar\_mass} — \texttt{rho\_critical} is
molar, $17\,873.7\ \mathrm{mol/m^3}$). Since Water's reducing temperature
and density \emph{are} its critical values, this makes $\delta = \tau = 1$
exactly — the branch point Part 2 discussed.

\section{Where Water's data comes from: \texttt{fluid.eos()}}

\texttt{fluid.eos()} (\texttt{src/fluid.rs}) is a \texttt{match} returning
\texttt{\&fluids::water::WATER} — a \texttt{const FluidEos} generated by
\texttt{dev/gen\_fluid.py} from CoolProp's own \texttt{Water.json}
(\texttt{upstream\_source/CoolProp/dev/fluids/Water.json}). Its residual
term list includes exactly one \texttt{ResidualTerm::NonAnalytic} entry,
carrying the two-term coefficient arrays from Part 1
(\texttt{a: \&[3.5, 3.5], b: \&[0.85, 0.95], beta: \&[0.3, 0.3], big\_a:
\&[0.32, 0.32], big\_b: \&[0.2, 0.2], big\_c: \&[28.0, 32.0], big\_d:
\&[700.0, 800.0]}) verbatim from CoolProp's fitted values.

\section{\texttt{eos.residual\_derivs}: summing every term type}

\texttt{FluidEos::residual\_derivs(delta, tau)} (\texttt{src/eos.rs}) walks
\emph{every} entry in \texttt{self.residual} and calls
\texttt{term.accumulate(delta, tau, \&mut acc)} on each — for Water this
means one \texttt{Power}-variant call (51 monomial terms), one
\texttt{Gaussian}-variant call (3 bell-shaped terms), and the
\texttt{NonAnalytic}-variant call from Part 3. All three add into the same
accumulator; \texttt{residual\_derivs} has no idea, and does not need to
know, which term types actually contributed.

\section{Inside the non-analytic call}

With $\delta=\tau=1$ exactly, \texttt{accumulate\_non\_analytic}'s guard
(Part 2) fires on both variables, offsetting each by $10\varepsilon \approx
2.22\times10^{-15}$. For term $i=1$ ($a_1=3.5$, $b_1=0.85$,
$\beta_1=0.3$, \ldots): $\mathrm{dm1} = \delta - 1 \approx 2.22\times
10^{-15}$, so $\mathrm{dm1\_sq} \approx 4.9\times10^{-30}$ — an extremely
small but perfectly finite \texttt{f64}. Propagating this through $\theta$,
$\Delta$, $\Delta^{b_i}$ and $\psi$ (Part 2's chain) and their derivatives
gives six finite numbers per term; both terms' contributions are summed into
\texttt{acc} via \texttt{add\_scaled}.

\section{Back up the stack: assembling the pressure}

Once \texttt{residual\_derivs} returns, \texttt{state\_trho} combines the
residual $\alpha_\delta$ (called \texttt{res.ad} in code) with the ideal-gas
part via the standard Span--Wagner relation
\[
  p = \rho\, R\, T\, (1 + \delta\, \alpha_{r,\delta}),
\]
implemented directly as \texttt{rho\_molar * r * t * (1.0 + delta * res.ad)}.

\section{The verified result}

\texttt{tests/non\_analytic\_critical\_region.rs} runs exactly this
computation and asserts the result against IAPWS-95's defining critical
pressure:
\begin{lstlisting}[language=Rust]
let s = state_trho(Fluid::Water, eos.t_critical, rho_mass_c);
let rel = (s.pressure - eos.p_critical).abs() / eos.p_critical;
assert!(rel < 1e-9);
\end{lstlisting}
Running it (\texttt{cargo test --release -p outram-park-fork-coolprop
--test non\_analytic\_critical\_region -- --nocapture}) prints:
\begin{center}
\texttt{p(Tc, rho\_c) = 22064000.00000 Pa (p\_c = 22064000 Pa), rel err = 5.217e-14}
\end{center}
five orders of magnitude tighter than the $10^{-9}$ gate, and — as
documented in this crate's git history — a substantial improvement over the
$\sim 10^{-4}$ residual the same computation carried before
\texttt{accumulate\_non\_analytic} was implemented (when the
\texttt{NonAnalytic} match arm was a documented no-op). The full
verification record, including the reference BibTeX citation with page
number, lives in
\texttt{verification\_and\_validation/water\_critical\_point\_iapws95.md}.

\end{document}
